% IV. Методика и експериментални резултати.
% Слети бивши раздели V (методика) и VI (резултати).
% Всяко число тук идва от изпълнение. Необработеният изход е в docs/measurements/.
\section{Методика и експериментални резултати}
\label{sec:results}

\subsection{Какво се мери и как}

Експериментът проверява три твърдения: че системата превръща вярно въпросите във формална
заявка, че логическата форма е действителна абстракция, а не преоблечен шаблон над
въпроса, и че перифразата съответства на построената заявка.

Наборът от 25 въпроса е замразен преди каквото и да е измерване и се чете както от
демонстрацията, така и от програмата за мерене. Той съдържа 5 въпроса в активен залог към
допълнението, 4 в активен залог към подлога, 3 в страдателен залог, по 2 с отрицание, с
броене, с отговор да или не, с квантор и с предложно допълнение, и 3 нарочно извън обхвата,
за да има покритието върху какво да се провали. Ресурсите са 56 лексеми, 80 словоформи, 25
продукции, 6 релации, 8 предложни рамки и 38 факта.

Мерките са пет: покритие на синтактичния анализ, покритие от край до край, нееднозначност
(изводи, разчитания след типовата проверка и различни логически форми, отчетени поотделно,
защото разликите между тях са приносът на семантичния слой), съгласуваност на генерирания
SQL и време за анализ. Времената се мерят със \code{std::chrono::steady\_clock} около
самото извикване на анализатора, без еднократното зареждане на ресурсите, и се повтарят
нечетен брой пъти с отчитане на медианата, защото еднократно измерване на анализ под една
милисекунда мери часовника, а не анализатора.

Числата са получени на 19 август 2026 г. при превод с g++ 15.2.0 (MinGW-w64) в режим
\code{Release} по стандарт C++20, CMake 4.3.2 и Ninja 1.13.2 върху Windows 11 с 12
логически процесора, а проверката на SQL с модула \code{sqlite3} на Python 3.14.
Необработеният изход е в \code{docs/measurements/} и е част от хранилището, тоест всяко
число се проследява до реда, от който идва.

\subsection{Работата на конвейера върху един въпрос}

Листинг~\ref{lst:trace} показва изхода за нееднозначен въпрос, съкратен до съществените
стъпки. Държавата може да принадлежи на актьорите или на филмите, моделът има релация и за
двете, така че и двете разчитания оцеляват и дават различен отговор. Системата избира
по-близкото прикрепване и казва кое и защо.

\begin{lstlisting}[language={},numbers=none,basicstyle=\scriptsize\ttfamily,
caption={Съкратен изход на \code{lexicon} за нееднозначен въпрос},label={lst:trace}]
QUESTION  Which films with actors from France did Nolan direct?

2. PARSE FOREST
   derivations 2   forest nodes 28   chart items 111   ambiguity points 1

3. READINGS
   reading 1: kept, attachment cost 4
      select x0 where Film(x0) & Person(x1) & acted_in(x1, x0)
                    & film_country(x0, c_france) & directed(p_nolan, x0)
   reading 2: kept, attachment cost 2
      select x0 where Film(x0) & Person(x1) & person_country(x1, c_france)
                    & acted_in(x1, x0) & directed(p_nolan, x0)

4. CHOSEN READING
   chose reading 2 with attachment cost 2, the closest attachment among
   the surviving readings

6. SQL
   SELECT DISTINCT t0.title AS answer
   FROM film AS t0, person AS t1, country AS e2, person AS e3,
        acted_in AS r4, directed AS r5
   WHERE e2.id = 'c_france' AND e3.id = 'p_nolan'
     AND t1.country_id = e2.id
     AND r4.person_id = t1.id AND r4.film_id = t0.id
     AND r5.person_id = e3.id AND r5.film_id = t0.id  ORDER BY 1;

7. SPARQL
   SELECT DISTINCT ?x0 WHERE { ?x0 a :Film . ?x1 a :Person .
     ?x1 :comesFrom :c_france . ?x1 :actedIn ?x0 .
     :p_nolan :directed ?x0 . } ORDER BY ?x0

8. PARAPHRASE BACK TO ENGLISH
   You are asking for the films with actors from France that were directed
   by Christopher Nolan.

9. ANSWER FROM THE BUNDLED KNOWLEDGE BASE
   Inception
\end{lstlisting}

\subsection{Покритие и съгласуваност по слоеве}

Разбивката по слоеве, Табл.~\ref{tab:bylayer}, отнася грешката към конкретен слой, вместо
да я остави обща за системата.

\begin{table}[H]
\caption{Покритие и съгласуваност по слоеве}
\label{tab:bylayer}
\centering
\begin{tabular}{@{}L{5.4cm}L{4.6cm}r@{}}
\toprule
\textbf{Слой} & \textbf{Мярка} & \textbf{Резултат} \\
\midrule
морфология & въпроси, в които всяка дума е позната & 22 от 25 \\
синтактичен анализ & въпроси с поне един извод & 23 от 25 \\
семантично изображение & въпроси с поне едно разчитане върху модела & 22 от 25 \\
генериране на SQL & заявки, които връщат същото като оценителя & 22 от 22 \\
генериране на SPARQL & произведена заявка за всеки построен израз & 22 от 22 \\
синтез & произведена перифраза за всеки построен израз & 22 от 22 \\
\bottomrule
\end{tabular}
\end{table}

Съгласуваността на SQL е измерена независимо: заявката се изпълнява върху SQLite, зареден
със схемата и данните от \code{lexicon --emit-sql-schema}, а резултатът се сравнява с
отговора на оценителя, който работи направо върху логическата форма и не вижда SQL. 22 от
22 съвпадат, включително отрицанията и релационното деление, което идва от \term{every}.
Всичките 30 теста минават. Трите въпроса с непозната дума са същите три, които са нарочно
извън обхвата. Отделна мярка за точност на морфологичния анализ спрямо ръчно проверени
разчитания \textbf{не е провеждана}: редът в Табл.~\ref{tab:bylayer} отчита само дали
думите са познати, а не дали приписаните признаци са верни за произволен текст.

\subsection{Поведение при нееднозначност}

Шест въпроса допускат повече от един извод (Табл.~\ref{tab:ambiguity}). При долните три
разчитанията се различават по дървото, но не по значението, и каноничното сравнение ги
събира в едно. При горните три разчитанията са наистина различни и всяко дава различен
отговор.

\begin{table}[H]
\caption{Разчитания преди и след семантичната проверка}
\label{tab:ambiguity}
\centering
\begin{tabular}{@{}L{7.4cm}rrr@{}}
\toprule
\textbf{Въпрос} & \textbf{изводи} & \textbf{върху модела} & \textbf{значения} \\
\midrule
Which actors acted in every film by Villeneuve? & 2 & 2 & 2 \\
Which actors acted in no film by Nolan? & 2 & 2 & 2 \\
Which films with actors from France did Nolan direct? & 2 & 2 & 2 \\
Which actors from France acted in films by Nolan? & 2 & 2 & 1 \\
Who directed a film with Cotillard in 2010? & 3 & 3 & 1 \\
Which directors from Canada directed films in 2016? & 2 & 2 & 1 \\
\bottomrule
\end{tabular}
\end{table}

Отхвърляне по типова проверка не се среща в замразения набор, но механизмът се проверява
от тестовете. Върху \term{Which films with actors in 2010 did Nolan direct?} синтактичният
слой дава 2 извода, а семантичният отхвърля единия с причината „no relation for in between
Person and Year“. Този въпрос не е част от набора и числото не влиза в
Табл.~\ref{tab:bylayer}.

\subsection{Нееднозначност спрямо размера на граматиката}

Вторият експеримент минава същия набор през три граматики. Втората добавя две продукции към
работните 25, третата още една, и никоя не описва низ, който първата не описва.

При работната граматика наборът дава 30 извода, 25 значения и 7 точки на нееднозначност.
Двете допълнителни продукции вдигат изводите на 47, тоест с 57 на сто, а точките на 19, без
да променят нито покритието (23 анализирани въпроса и при трите), нито броя на значенията.
Третата граматика не променя нищо: добавената продукция за две съседни съществителни не се
задейства, защото нито един въпрос не поставя две съществителни едно до друго. Цената на
една продукция следователно не е самото ѝ съществуване, а това дали изреченията стигат до
нея. Значенията остават 25 и при трите, но си струва да се каже без заобикалки: излишните
разчитания отпадат, защото за тези продукции няма написано семантично правило, а не защото
моделът на знанието ги е отхвърлил.

\subsection{Време за анализ}

Табл.~\ref{tab:parsetime} дава медианата и минимума от 201 повторения за изречения, всяко
от които добавя по една предложна група към предходното.

\begin{table}[H]
\caption{Време за анализ спрямо дължината на изречението, 201 повторения}
\label{tab:parsetime}
\centering
\begin{tabular}{@{}rrrrrrr@{}}
\toprule
\textbf{групи} & \textbf{думи} & \textbf{изводи} & \textbf{елементи} & \textbf{възли} & \textbf{медиана, us} & \textbf{минимум, us} \\
\midrule
0 &  4 & 1 &  48 & 13 & 14.6 & 14.1 \\
1 &  6 & 2 &  66 & 22 & 22.7 & 21.9 \\
2 &  8 & 3 &  84 & 31 & 30.0 & 26.5 \\
3 & 10 & 4 & 102 & 40 & 35.3 & 33.6 \\
4 & 12 & 5 & 120 & 49 & 45.4 & 42.3 \\
5 & 14 & 6 & 138 & 58 & 52.8 & 49.6 \\
6 & 16 & 7 & 156 & 67 & 64.7 & 55.9 \\
\bottomrule
\end{tabular}
\end{table}

Изводите растат като $k+1$ за $k$ предложни групи, а елементите и възлите линейно, с точно
18 елемента и 9 възела на добавена група, и времето върви заедно с тях. Кубичната горна
граница не се проявява в този диапазон, защото изреченията са къси, а граматиката не
поражда плътна таблица. Средното време за един въпрос от набора, измерено веднъж за всеки
въпрос, е 45,1 микросекунди; това число е по-шумно, защото не е повторено.

\subsection{Съпоставка с базово решение и анализ на провалите}

Съпоставката е проведена. Условието, което този раздел поставяше, е фиксиран модел,
фиксирана версия и фиксирани настройки, защото без тях резултатът не е възпроизводим.
Локален модел изпълнява това условие по-добре от услуга през чужд програмен интерфейс:
отдалечената крайна точка се сменя без предупреждение и старото поведение е
невъзстановимо, докато модел, изтеглен по контролна сума и пуснат при температура 0 с
фиксирано зърно, се повтаря. Затова базовото решение тук е локално.

Моделът е \code{qwen2.5-coder:7b}, контролна сума \code{dae161e27b0e}, при температура 0,
зърно 1 и \code{num\_predict} 400, пуснат през Ollama 0.32.9. Програмата за мерене е
\code{tools/llm\_baseline.py}, а второто условие се получава с ключа \code{--with-labels}.
Необработеният изход е в \code{docs/measurements/llm\_baseline\_schema\_only.txt} и
\code{docs/measurements/llm\_baseline\_with\_labels.txt}. Повторимостта е проверена, не допусната: двете условия
бяха пуснати два пъти и присъдите по всичките 25 въпроса съвпаднаха.

Критерият е равенство по изпълнение, а не по низ. Заявката на конвейера и заявката на
модела се изпълняват върху вградената база и двата резултата се сравняват като неподредени
множества с повторения. Две различни заявки, които връщат едни и същи редове, се броят и
двете за верни; сравнение по низ би отчело верни отговори като грешни. Условията са две
и разликата между тях е основният резултат (Табл.~\ref{tab:llmbaseline}).

\begin{table}[H]
\caption{Съвпадения на локалния езиков модел с конвейера, 25 въпроса}
\label{tab:llmbaseline}
\centering
\begin{tabular}{@{}L{7.4cm}rr@{}}
\toprule
\textbf{Условие} & \textbf{съвпадения} & \textbf{неизпълними заявки} \\
\midrule
само схемата & 5 от 25 & 1 \\
схемата и стойностите на етикетите в данните & 16 от 25 & 2 \\
\bottomrule
\end{tabular}
\end{table}

Второто условие съществува, защото правиловият конвейер получава лексикон, който свързва
повърхнинните форми с обектите. Да се отнемат на модела всички етикети, докато другата
страна има лексикон, не е съпоставка при равни условия. Подаването на етикетите утроява
съвпаденията, от 5 на 16. Това показва къде е бил основният провал: в свързването на
имената с обектите, а не в построяването на заявката. Оттам се вижда и къде стои
действителната стойност на лексикона в тази работа.

Деветте разлики при подадени етикети са изброени поименно (Табл.~\ref{tab:llmdiffs}).
Разбивката казва повече от самото число и на две места говори срещу проекта.

\begin{table}[H]
\caption{Деветте разлики при подадени етикети}
\label{tab:llmdiffs}
\centering
\begin{tabular}{@{}L{6.0cm}L{5.6cm}L{2.6cm}@{}}
\toprule
\textbf{Въпрос} & \textbf{Какво се разминава} & \textbf{Причина} \\
\midrule
Did Nolan direct Inception? & 1 срещу „Yes“ & критерият \\
Did Villeneuve direct Inception? & 0 срещу „No“ & критерият \\
Which films did Nolan direct or produce? & конвейерът не дава заявка, моделът е верен & конвейерът \\
Which films that Nolan directed were released in 2010? & конвейерът не дава заявка, моделът е верен & конвейерът \\
Which films did Kubrick direct? & конвейерът не дава заявка, моделът връща празно, което е вярно & конвейерът \\
Who directed Inception? & измислена колона \code{person.person\_id} & моделът \\
Which directors from Canada directed films in 2016? & погрешно назована \code{T3.name} & моделът \\
Which actors from France acted in films by Nolan? & празно вместо Marion Cotillard & моделът \\
Which actors acted in every film by Villeneuve? & режисьорът вместо актьора & моделът \\
\bottomrule
\end{tabular}
\end{table}

Два от деветте не са провал на модела. При \term{Did Nolan direct Inception?} и \term{Did
Villeneuve direct Inception?} конвейерът връща целите числа 1 и 0, а моделът низовете
„Yes“ и „No“. Значението е едно и също, но множествата от редове се различават. Това е
ограничение на критерия и се записва като такова; критерият не се разхлабва след като
числата са известни.

Три от деветте са там, където конвейерът изобщо не произвежда SQL. Моделът отговаря вярно
на \term{Which films did Nolan direct or produce?} и \term{Which films that Nolan directed
were released in 2010?}, а на \term{Which films did Kubrick direct?} връща празно
множество, което е вярно, защото Кубрик не е в данните. Тоест моделът успява точно там,
където правиловата система се проваля.

Четири са действителни провали на модела. При \term{Who directed Inception?} и \term{Which
directors from Canada directed films in 2016?} той измисля колони, които не съществуват, и
заявката не се изпълнява. При \term{Which actors from France acted in films by Nolan?}
връща празно множество, а отговорът е Marion Cotillard, при съединение през две стъпки с
филтър по държава. При \term{Which actors acted in every film by Villeneuve?} връща
режисьора вместо актьора, тоест се проваля на квантора; обхватът на кванторите вече стои в
по-нататъшната работа.

Тоест от 22-те въпроса, на които конвейерът отговаря, моделът съвпада по резултат на 16,
на още два дава семантично същия отговор в друг запис, и на четири се проваля
действително.

Границите се казват, а не се подразбират. Моделът е достатъчно малък, за да се побере в
6 GB видеопамет. Числата казват какво прави един малък квантуван локален модел върху 25
въпроса срещу ръчно написан каталог от филми, и нищо за челен модел или за езиковите
модели изобщо.

Провалите на конвейера са три и вече се появиха по име. Първите два се провалят в
синтактичния слой, защото граматиката няма съчинение и няма относително изречение, а
третият в семантичния, защото името не е в лексикона. Първите два се поправят с продукция
и семантично правило, третият с един ред данни. Нито един провал не идва от логическата
форма или от генераторите. Покритието на граматиката следователно е по-скъпото
ограничение, и всяко негово разширение носи нова нееднозначност с измерената по-горе цена.
Че точно на тези три въпроса базовото решение отговаря вярно, дава мащаб на тази цена.
